\begin{statementbox}
\textbf{C008\quad 椭圆切点式、切线与极线}

设椭圆
\[
C:\quad \frac{x^2}{a^2}+\frac{y^2}{b^2}=1,\qquad a>b>0.
\]

\begin{enumerate}
    \item \textbf{已知切点，直接写切线。}
    若 $P(x_0,y_0)\in C$，即
    \[
    \frac{x_0^2}{a^2}+\frac{y_0^2}{b^2}=1,
    \]
    则椭圆在 $P$ 处的切线为
    \[
    \boxed{\frac{x_0x}{a^2}+\frac{y_0y}{b^2}=1.}
    \]
    此式称为\textbf{切点式}：把椭圆方程中的 $x^2,y^2$ 分别换成 $x_0x,y_0y$，分母和右边的 $1$ 不变。

    \item \textbf{切线方向与斜率。}
    切线的一个法向量为
    \[
    \boldsymbol n=\left(\frac{x_0}{a^2},\frac{y_0}{b^2}\right),
    \]
    一个方向向量为
    \[
    \boldsymbol d=\left(-\frac{y_0}{b^2},\frac{x_0}{a^2}\right).
    \]
    若 $y_0\ne0$，切线斜率为
    \[
    \boxed{k=-\frac{b^2x_0}{a^2y_0}.}
    \]
    若 $y_0=0$，则 $P=(\pm a,0)$，切线为 $x=\pm a$；若 $x_0=0$，则 $P=(0,\pm b)$，切线为 $y=\pm b$。

    \item \textbf{已知椭圆外点，写两切点的连线。}
    设 $M(x_1,y_1)$ 满足
    \[
    \frac{x_1^2}{a^2}+\frac{y_1^2}{b^2}>1.
    \]
    从 $M$ 向椭圆作两条实切线，切点为 $T_1,T_2$，则\textbf{切点弦（极线）} $T_1T_2$ 的方程为
    \[
    \boxed{\frac{x_1x}{a^2}+\frac{y_1y}{b^2}=1.}
    \]
    注意：这条直线经过的是\textbf{两个切点}，一般\textbf{不经过外点 $M$}，也不是从 $M$ 引出的两条切线之一。

    \item \textbf{极线的完整位置分类。}
    对非零点 $M(x_1,y_1)$，记
    \[
    S=\frac{x_1^2}{a^2}+\frac{y_1^2}{b^2}.
    \]
    方程 $x_1x/a^2+y_1y/b^2=1$ 在 $S>1$ 时与椭圆交于两个实切点；在 $S=1$ 时就是 $M$ 处的切线；在 $0<S<1$ 时仍是一条实直线，但\textbf{不与椭圆相交，也不存在从 $M$ 引出的实切线}。若 $M=(0,0)$，代入得到 $0=1$，它\textbf{不是普通平面中的直线}。

    \item \textbf{拓展：已知直线，判断是否相切。}
    对 $L:Ax+By+C=0$，其中 $(A,B)\ne(0,0)$，有
    \[
    \boxed{L\text{ 与 }C\text{ 相切}\iff C^2=a^2A^2+b^2B^2.}
    \]
\end{enumerate}
\end{statementbox}
